This section states the planning deadlines of Section~\ref{sec:method} of the main paper in full and proves
Propositions~\ref{prop:cc} and~\ref{prop:costfb}. For task $j$, $f_j$ is the planned finish time, $d_j$ the deadline,
$w_j$ the reward, $h_j=f_j-t_0$ the lead time and $F_j=f_j+e_j$ the actual finish time, with
$e_j\sim\mathcal{N}(m(h_j),s(h_j)^2)$ and $m$, $s$ as in Eq.~\eqref{eq:errmodel}. The effective deadline of
Eq.~\eqref{eq:cc} at safety factor $z$ is written $\tilde d_j(z)$, and $\tilde d_j(-\infty)$ is read as $d_j$.

\subsection{Planning deadlines}
Let $\tilde H$ be the set of tasks that no vehicle can finish by $\tilde d_j(z)$ on a direct flight, and $H$ the set of
tasks that no vehicle can finish by $d_j$. The nominal fallback plans against
\begin{equation}\label{eq:Sfallback}
d'_j=\begin{cases} d_j, & j\in\tilde H\setminus H,\\ \tilde d_j(z), & \text{otherwise.}\end{cases}
\end{equation}
For the cost-aware fallback, let $\theta_j=c_j/(w_j+c_j)$ be the break-even probability of Eq.~\eqref{eq:theta},
$z_j=\Phi^{-1}(\theta_j)$ with $z_j=-\infty$ for $c_j=0$, and $z_{p,j}=\max(z,z_j)$. With $\tilde H_p$ the set of tasks
that no vehicle can finish by $\tilde d_j(z_{p,j})$ on a direct flight, the planning deadline is
\begin{equation}\label{eq:Scostfb}
d'_j=\begin{cases} \min\{d_j,\tilde d_j(z_j)\}, & j\in\tilde H_p\setminus H,\\ \tilde d_j(z_{p,j}), &
\text{otherwise.}\end{cases}
\end{equation}
A task that cannot meet this deadline either is given up. Finding $\tilde H_p$ and $H$ requires one direct-flight check
per task and vehicle. Where $1+b_1+z_js_1\le0.05$, the effective deadline is not defined and the nominal deadline is
used, without the guarantee of part (iii) of Proposition~\ref{prop:costfb}. With the error models of this paper, this
happens only for $\theta_j<0.03$, which none of the studied late costs produces. Every planner caps the effective deadline
at $d_j$, so that no task is planned beyond its nominal deadline.

\subsection{Proposition 1}
In full, Proposition~\ref{prop:cc} states the following. Let $1+b_1+zs_1>0$ and $s_0,s_1\ge0$. (i) A finish time
satisfies $f_j+m(f_j-t_0)+z\,s(f_j-t_0)\le d_j$ if and only if $f_j\le\tilde d_j(z)$. (ii) Under the error model, a task
planned to finish at $f_j\le\tilde d_j(z)$ meets its deadline with probability at least $\Phi(z)$ if its route starts as
planned. (iii) With $b_0=b_1=s_0=s_1=0$, $\tilde d_j(z)=d_j$ and the planning model is the nominal one.

\begin{proof}
Substituting the error model gives $f_j(1+b_1+zs_1)\le d_j-b_0-zs_0+(b_1+zs_1)t_0$, and dividing by the positive factor
gives (i). For (ii), $F_j=f_j+e_j$ with $e_j\sim\mathcal{N}(m(h_j),s(h_j)^2)$. By (i), $d_j-f_j-m(h_j)\ge z\,s(h_j)$. If
$s(h_j)>0$, then $\Pr[F_j\le d_j]=\Phi\big((d_j-f_j-m(h_j))/s(h_j)\big)\ge\Phi(z)$. If $s(h_j)=0$, which happens only
for $s_0=0$ and $h_j=0$, the error equals $m(h_j)$ with certainty, $F_j\le d_j$ holds, and the probability is one. Part
(iii) is immediate.
\end{proof}

The cap at $d_j$ matters only where $\tilde d_j(z)>d_j$, which happens when $z<0$ or the mean error is negative. It then
adds the nominal deadline to the constraint, which becomes tighter, so part (ii) still holds, whereas the equivalence of
part (i) refers to the uncapped deadline.

\subsection{Proposition 2}
In full, Proposition~\ref{prop:costfb} states the following. Let $s_0,s_1\ge0$, $1+b_1+zs_1>0$ and $1+b_1+z_js_1>0$
for every finite $z_j$. (i) If $c_j=0$ for every task, Eq.~\eqref{eq:Scostfb} is the nominal fallback of
Eq.~\eqref{eq:Sfallback}. (ii) A task with $\theta_j\ge\Phi(z)$ never falls back, and it is planned against
$\tilde d_j(z_j)$. (iii) Under the error model, every task planned to finish by $d'_j$ has an expected contribution
$w_j\Pr[F_j\le d_j]-c_j\Pr[F_j>d_j]\ge0$ if its route starts as planned. (iv) If every $c_j$ is non-decreasing in
$\lambda$, every planning deadline $d'_j$ is non-increasing in $\lambda$.

\begin{proof}
(i) $z_j=-\infty$ gives $z_{p,j}=z$ and $\min\{d_j,\tilde d_j(-\infty)\}=d_j$. (ii) $z_j\ge z$ holds if and only if
$\theta_j\ge\Phi(z)$. Then $z_{p,j}=z_j$, and a task that cannot meet $\tilde d_j(z_{p,j})$ cannot meet
$\min\{d_j,\tilde d_j(z_j)\}$ either, so it is given up. (iii) A protected task meets its deadline with probability at
least $\Phi(z_{p,j})\ge\Phi(z_j)$, and a fallback task with probability at least $\Phi(z_j)$, both by
Proposition~\ref{prop:cc}, whose part (ii) also holds with the cap at $d_j$. Hence $p\ge\theta_j$, and the expected
contribution $(w_j+c_j)p-c_j$ is not negative. (iv) $\theta_j$ and $z_j$ do not decrease with $\lambda$, and with
$s_0,s_1\ge0$, $\tilde d_j(z)$ does not increase with $z$. A fallback task stays one when $\lambda$ grows, because
$z_{p,j}$ does not decrease. A protected task can become a fallback task only while $z_j<z$, and then $z_{p,j}=z$ does not
change. In each case the deadline does not increase.
\end{proof}

Part (iii) is a local criterion. It prices each attempt on its own and ignores the delay that a fallback task causes to
the tasks after it on the same route, as well as the reassignments it may trigger, so it does not guarantee that
attempting the task raises the utility of the whole plan.
